% Part:sets-functions-relations
% Chapter: size-of-sets
% Section: enumerations-alt

\documentclass[../../../include/open-logic-section]{subfiles}

\begin{document}

\olfileid{sfr}{siz}{enm-alt}

\olsection{شمېرنې او \usetoken{S}{enumerable} سټونه}

\begin{editorial}
  دا برخه شمېرنې له $\Nat$ (پيلنيو ټوټو) سره د دوه اړخيزو يو
  پر يو تابعو په توګه تعريفوي، لکه په سټ‌تيورۍ کښې چې کېږي۔ نو
  د \olref[enm]{sec} له تعريفونو سره لږ ټکر لري، او هلته ټولې
  بېلګې تکراروي۔ دا د هغې برخې په پرتله لږه لنډه هم ده۔
\end{editorial}

يو متناهي سټ يوازې د هغه د !!{element}s په شمېرلو ټاکلے شو۔ دا
هغه وخت کوو چې يو سټ داسې تعريفوو:
\[
  A = \{a_1, a_2, \ldots, a_n\}.
\]
که فرض کړو چې !!{element}s $a_1$، \dots، $a_n$ ټول سره بېل دي،
دا د $A$ او لومړنيو $n$ طبيعي عددونو $0$، \dots، $n-1$ تر منځ
يوه !!a{bijection} راکوي۔ برعکس، ځکه هر متناهي سټ يوازې متناهي
شمېر !!{element}s لري، هر متناهي سټ په داسې سمون کښې اېښودلے
شو۔ په بله وينا، که $A$ متناهي وي، د $A$ او
$\{0, \dots, n-1\}$ تر منځ يوه !!a{bijection} شته، چې $n$ د~$A$
د !!{element}s شمېر دے۔

که د نامتناهي سټونو ځينې ډولونه ومنو، نو ځينې نامتناهي سټونه به
هم شمېرل کېدے شي۔ دا په دې وينا دقيقولے شو چې يو نامتناهي سټ د
هغه او ټول~$\Nat$ تر منځ د !!a{bijection} په وسيله شمېرل کېږي۔

\begin{defn}[شمېرنه، سټ‌تيوريکي]
د يوه سټ $A$ \emph{شمېرنه} يوه !!a{bijection} ده چې د قيمتونو
سټ ئې $A$ او د تعريف ساحه ئې يا د طبيعي عددونو يو پيلنی سټ
$\{0, 1, \ldots, n\}$ {يا} د طبيعي عددونو ټول سټ~$\Nat$ وي۔
\end{defn}

\begin{explain}
د \emph{شمېرنې} د ويي د دې کارونې تر شا يو شهودي بنسټ شته۔ دا
ويل چې موږ يو سټ $A$ شمېرلے دے، په دې معنا دي چې يوه
!!a{bijection} $f$ شته چې پرې د سټ $A$ غړي شمېرلے شو۔ $0$م غړے
$f(0)$ دے، لومړے $f(1)$، \ldots، $n$م $f(n)$،
\ldots۔\footnote{هو، موږ له $0$ څخه شمېرو۔ البته له~$1$ څخه هم
پيلولے شو۔ په دې سره به لوے توپير رانۀ شي۔ يوازې به مو~$\Nat$
په~$\PosInt$ بدلولے۔} د دې دليل د لاندې خبرې په ورزياتولو لا
روښانه کېدے شي:
\end{explain}

\begin{defn}
  \ollabel{defn:enumerable}
  يو سټ~$A$ هله او يوازې هله !!{enumerable} دے چې يا
  $A = \emptyset$ وي يا د~$A$ يوه شمېرنه وي۔ وايو چې $A$ هله
  او يوازې هله !!{nonenumerable} دے چې $A$ !!{enumerable} نۀ وي۔
\end{defn}

\begin{explain}
نو يو سټ هله او يوازې هله !!{enumerable} دے چې تش وي يا ئې
!!{element}s په يوه شمېرنه شمېرلے شئ۔
\end{explain}

\begin{ex}
هغه تابع چې طبيعي عددونه شمېري يوازې عيني تابع
$\Id{\Nat} \colon \Nat \to \Nat$ ده چې په
$\Id{\Nat}(n) = n$ ورکړل شوې۔ هغه تابع چې \emph{مثبت} طبيعي
عددونه شمېري، $\Nat^+ = \Nat \setminus \{0\}$، تابع
$g(n) = n + 1$ ده؛ يعنې راتلونکې تابع۔
\end{ex}

\begin{prob}
وښيئ چې يو سټ $A$ هله او يوازې هله !!{enumerable} دے چې يا
$A = \emptyset$ وي يا يوه !!a{surjection} $f\colon \Nat \to A$
وي۔ وښيئ چې $A$ هله او يوازې هله !!{enumerable} دے چې يوه
!!a{injection} $g\colon A \to \Nat$ وي۔
\end{prob}

\begin{ex}
تابعې $f\colon \Nat \to \Nat$ او $g \colon \Nat \to \Nat$ چې
په ترتيب سره په لاندې ډول ورکړل شوې دي:
\begin{align*}
  f(n) & = 2n \text{ او}\\
  g(n) & = 2n+1
\end{align*}
جفت طبيعي عددونه او طاق طبيعي عددونه شمېري۔ خو له دواړو يوه هم
!!{surjective} نۀ ده، نو يوه هم د $\Nat$ شمېرنه نۀ ده۔
\end{ex}

\begin{prob}
د مربع عددونو $1$، $4$، $9$، $16$، \dots يوه شمېرنه تعريف کړئ۔
\end{prob}

\begin{ex}
فرض کړئ $\lceil x \rceil$ د \emph{پاسني صحيح عدد} تابع ده، چې
$x$ تر ټولو نږدې صحيح عدد ته پورته ګردوي۔ بيا تابع
$f \colon \Nat \to \Int$ چې په لاندې ډول ورکړل شوې ده:
\[
  f(n) = (-1)^{n} \left\lceil\tfrac{n}{2}\right\rceil
\]
د صحيح عددونو سټ~$\Int$ په لاندې ډول شمېري:
\[
\begin{array}{c c c c c c c c}
f(0) & f(1) & f(2) & f(3) & f(4) & f(5) & f(6) & \dots \\ \\
\big\lceil \tfrac{0}{2} \big\rceil & -\big\lceil \tfrac{1}{2}\big\rceil &  \big\lceil \tfrac{2}{2} \big\rceil & -\big\lceil \tfrac{3}{2} \big\rceil & \big\lceil \tfrac{4}{2} \big\rceil  & -\big\lceil \tfrac{5}{2}\big\rceil & \big\lceil \tfrac{6}{2} \big\rceil & \dots \\ \\
0 & -1 & 1 & -2 & 2 & -3 & 3& \dots
\end{array}
\]
پام وکړئ چې $f$ څنګه د مثبتو او منفي صحيح عددونو تر منځ په
«ټوپونو» د~$\Int$ قيمتونه پيدا کوي۔ $f$ په حالتونو هم داسې
تعريفولے شئ:
\[
f(n) = \begin{cases}
  \frac{n}{2} & \text{که $n$ جفت وي}\\
  -\frac{n+1}{2} & \text{که $n$ طاق وي}
  \end{cases}
\]
\end{ex}

\begin{prob}
وښيئ چې که $A$ او $B$ !!{enumerable} وي، نو $A \cup B$ هم دے۔
\end{prob}

\begin{prob}
پر $n$ په استقرا وښيئ چې که $A_1$، $A_2$، \dots، $A_n$ ټول
!!{enumerable} وي، نو $A_1 \cup \dots \cup A_n$ هم دے۔
\end{prob}

\end{document}
